\documentclass[10pt,a4paper]{amsart}
\usepackage[margin=19mm]{geometry}
\usepackage{amsmath,amssymb,mathtools}
\usepackage[scr=rsfs,cal=euler]{mathalfa}
\usepackage{xfrac}
\usepackage[unicode,hidelinks,bookmarks=false]{hyperref}
\newcommand{\ie}{i.e.\ }
\newcommand{\cf}{cf.\ }
\newcommand{\resp}{respectively\ }
\newcommand{\Subsection}{\subsection}
\providecommand{\nobreakdash}{\nobreak\hbox{-}}
\setlength{\emergencystretch}{2em}
\multlinegap0pt
\usepackage{amssymb}
\usepackage{mathtools}
\usepackage[shortlabels]{enumitem}
\newtheorem{proposition}{Proposition}
\def\b{\beta}
\newcommand{\mc}{\mathcal}
\def\o{\omega}
\def\s{\sigma}
\title{\textbf{Limiting behavior of principal eigenvalues for a class of mixed boundary value problems as the measure of the support domain goes to zero }}
\author{J. Lopez-Gomez, A. Sahuquillo}
\date{Source: arXiv:2603.17721v1 (CC BY 4.0)}
\begin{document}
\maketitle

\noindent\emph{Source and licence.} \textbf{Limiting behavior of principal eigenvalues for a class of mixed boundary value problems as the measure of the support domain goes to zero }. Version arXiv:2603.17721v1 (CC BY 4.0), distributed by arXiv under the Creative Commons Attribution 4.0 International licence, http://creativecommons.org/licenses/by/4.0/, as recorded in the arXiv licence metadata for this exact version and shown on the abstract page https://arxiv.org/abs/2603.17721v1. Full licence notice: \texttt{licenses/arxiv-2603.17721-CC-BY-4.0.txt}.\par\smallskip

\noindent\textit{Editorial scope.} Counted unit: the article's proposition pr2.3 (source0.tex lines 446--452). The article's complete original proof of it is delivered with it (source0.tex lines 453--475, one \texttt{proof} environment of 23 lines). No mathematics is written, repaired or re-derived: the statement and the proof are the article's own lines, in the article's own notation, and no retained line is substituted or re-flowed.  Retained uncounted as the source-local context the proof needs: lines 211--216.  Nothing was omitted from any retained span, so the package carries no omission record.  No retained line contains a citation, so the wrapper carries no bibliography and no bibliography entry.  No retained line contains a figure environment, an image-inclusion command or drawing code, so no image is delivered and the package declares no asset dependency.  Editorial formatting only, in addition: an AMS \texttt{amsart} preamble that restates the article's own declarations and macros used by the retained lines, namely the environments \texttt{proposition} on separate counters, together with the standard packages those lines need.  The retained environments print in this document as Proposition 1 in order of appearance, whereas the article prints the same items with the numbering of its own full document.  The complete native source of the article as distributed in the arXiv e-print for this exact version is bundled unchanged as \texttt{sources/196642/source0.tex}.
\begin{equation}
\label{1.3}
	\left \{ \begin{array}{ll}
		-u''=\s u \qquad  \hbox{in} \;\; [0,L],\\[1ex]
		\mc{B}_0 u(0) = 0, \quad \mc{B}_L u(L) = 0, \end{array} \right.
\end{equation}

\begin{proposition}
\label{pr2.3}
The principal eigenvalue of \eqref{1.3} is symmetric with respect to $\b_0\in(-\infty,+\infty]$ and $\b_\o\in(-\infty,+\infty]$, in the sense that
$$
\s_1[-D^2,\b_0,\b_\o,(0,L)] = \s_1[-D^2,\b_\o,\b_0,(0,L)].
$$
\end{proposition}

\begin{proof}
Let $u_1 \gg 0$ be a principal eigenfunction of \eqref{1.3}. Then, the change of variable
\begin{equation}
\label{2.2}
 v_1(y):= u_1(x) \quad \hbox{with} \; \; y := L-x, \; \; y \in [0,L],
\end{equation}
yields $u_1(0) = v_1(L)$ and $u_1(L) = v_1(0)$. Moreover, $u_1'(x) = -v_1'(y)$ and $u_1''(x) = v_1''(y)$ for all $x\in [0,L]$. In particular,
$$
  u_1'(0) = -v_1'(L),\quad u_1'(L) = -v_1'(0).
$$
Thus, the problem \eqref{1.3} becomes
\begin{equation}
\label{2.3}
	\left \{ \begin{array}{ll}
		-v_1''=\s_1[-D^2, \b_0 ,\b_\o,(0,L)] v_1  \qquad  \hbox{in} \;\; [0,L],\\[1ex]
		-v_1'(0)+\b_\o v_1(0)= 0, \quad  v_1'(L)+\b_0 v_1(L)  = 0. \end{array} \right.
\end{equation}
Since $u_1 \gg 0$, \eqref{2.2} implies that $v_1 \gg 0$. Hence, $v_1$ is a principal eigenfunction of \eqref{2.3} and, consequently, the uniqueness of the principal eigenvalue entails that
$$
\s_1[-D^2,\b_\o,\b_0,(0,L)] = \s_1[-D^2,\b_0,\b_\o,(0,L)].
$$
The proof is complete.
\end{proof}

\end{document}
